\section{Controlled Testbed and Experimental Design}
\label{sec:experiment_setup}

We study directional capability loss across three software engineering settings. 
The main
experiments use obfuscation and deobfuscation, which give controlled access to
forward pairs, executable validation, and inverse tasks of graded structural
complexity. 
We then evaluate whether the same phenomenon generalizes to two additional settings: Python-to-Java translation with Java-to-Python reversal, and C-to-assembly compilation with assembly-to-C reversal.
% Two further settings, Python-to-Java translation with Java-to-Python
% reversal and C-to-assembly compilation with assembly-to-C reversal, extend the
% analysis to inverse tasks of a different nature. Together they pair one detailed
% controlled study with two independent checks.

\subsection{Models}
\label{sec:models}

We evaluate the six LLMs in \Cref{tab:models}. Four are open-source 7B-class
models that support LoRA fine-tuning on academic hardware, and two are
proprietary instruction-tuned models that exposed a fine-tuning interface at the
time of the experiments. The set is not a leaderboard. It is chosen so that the
phenomenon can be tested across open and proprietary systems, general and
code-specialized pretraining, and reasoning-distilled and standard decoder-only
models. More recent and more capable models exist, but limited fine-tuning access and their size prevented us from using them; Section~\ref{sec:models-ttv} discusses this limitation.


\begin{table}[h]
\centering
\caption{Evaluated LLMs.}
\label{tab:models}
\begin{adjustbox}{width=0.9\columnwidth,center}
\begin{tabular}{@{}llll@{}}
\toprule
%\thead{Model} & \thead{Origin} & \thead{Specialization} & \thead{Fine-Tuning\\Technique} \\ \midrule
\textbf{Model} & \textbf{Origin} & \textbf{Specialization} & \textbf{Fine-Tuning} \\ \midrule
Qwen2.5-7B~\cite{qwen2} & Open-source & General-purpose & LoRA \\
Qwen2.5-Coder-7B~\cite{qwen2coder} & Open-source & Code-specialized & LoRA \\
DeepSeek-R1-Distill-Qwen-7B~\cite{deepseek} & Open-source & Reasoning-distilled & LoRA \\
Mistral-7B~\cite{mistral} & Open-source & General-purpose & LoRA \\
GPT-3.5-Turbo~\cite{openai2023gpt35} & Proprietary & Instruction-tuned & API \\
GPT-4.1-mini~\cite{openai2025gpt41mini} & Proprietary & Instruction-tuned & API \\ \bottomrule
\end{tabular}
\end{adjustbox}
\end{table}

\subsection{Obfuscation Tasks}
\label{sec:framework_instantiation}

The forward direction maps $C_{\text{orig}}$ to an obfuscated program
$C_{\text{trans}}$. The reverse direction asks whether a fine-tuned model can
map $C_{\text{trans}}$ back to a semantically equivalent program that is closer
to the original readable code. We implement three semantics-preserving
transformations with the open-source \emph{Obfuscator}
tool~\cite{superblaubeere27obfuscator}. Identifier Renaming is a lexical task
that replaces identifiers with obfuscated names. Dead Code Insertion is a
structural task that injects non-functional conditional blocks. Literal
Encryption is a semantic task.
It replaces literals with encrypted forms and embeds hard-coded decryption
logic, typically AES or an equivalent cipher in the Obfuscator tool.
A model cannot recover the plaintext by memorizing a general inverse mapping
alone. It must reason about the encrypted values and the embedded decryption
behavior.
% Encryption is a task that replaces literals with encrypted forms and
% hard-coded decryption logic. 
These techniques are common in practice.
Identifier renaming is the most widespread in Android apps~\cite{dong2018android}, and string encryption is frequent in malware~\cite{dong2018android}.
Junk code insertion is a standard strategy of Android obfuscators~\cite{hammad2018obfuscation}.
Figure~\ref{fig:transformation_examples} %\js{dead ref}
illustrates representative examples. Appendix provides the full taxonomy.

\begin{figure*}[t]
\centering
\lstset{language=Java, basicstyle=\scriptsize\ttfamily, numbers=none, xleftmargin=0pt, breaklines=true, aboveskip=2pt, belowskip=2pt}
\newcommand{\exlabel}[1]{\par\smallskip{\footnotesize\textbf{#1}}\par}
\newcommand{\exsep}{\hfill\rule[-6.0cm]{0.4pt}{6.0cm}\hfill}

\begin{subfigure}[t]{0.31\textwidth}
    \caption{Identifier Renaming}
    \label{fig:var_renaming}
    \exlabel{Original}
\begin{lstlisting}
int calculateSum(int[] numbers) {
    int total = 0;
    for (int num : numbers) {
        total += num;
    }
    return total;
}
\end{lstlisting}
    \exlabel{Obfuscated}
\begin{lstlisting}
int lIii(int[] Iiii) {
    int liii = 0;
    for (int iIIi : Iiii) {
        liii += iIIi;
    }
    return liii;
}
\end{lstlisting}
\end{subfigure}
\exsep
\begin{subfigure}[t]{0.31\textwidth}
    \caption{Dead Code Insertion}
    \label{fig:dead_code}
    \exlabel{Original}
\begin{lstlisting}
int getMax(int a, int b) {
    if (a > b) {
        return a;
    }
    return b;
}
\end{lstlisting}
    \exlabel{Obfuscated}
\begin{lstlisting}
int getMax(int a, int b) {
    if (false) {
        System.out.println("dead");
    }
    if (a > b) {
        return a;
    }
    return b;
}
\end{lstlisting}
\end{subfigure}
\exsep
\begin{subfigure}[t]{0.31\textwidth}
    \caption{Literal Encryption}
    \label{fig:literal_enc}
    \exlabel{Original}
\begin{lstlisting}
String greet(String name) {
    String msg = "Hello, ";
    if (name.isEmpty()) {
        name = "world";
    }
    return msg + name;
}
\end{lstlisting}
    \exlabel{Obfuscated}
\begin{lstlisting}
String greet(String name) {
    String msg = decrypt("Xy9K");
    if (name.isEmpty()) {
        name = decrypt("Zt8M");
    }
    return msg + name;
}
\end{lstlisting}
\end{subfigure}

\caption{Code transformation examples. Original and obfuscated snippets for Identifier Renaming, Dead Code Insertion, and Literal Encryption.}
\label{fig:transformation_examples}
\end{figure*}


The ordering from Identifier Renaming through Dead Code Insertion to Literal
Encryption underpins our interpretation. Identifier Renaming admits many
acceptable reversals, yet its inverse stays local and syntactic. Dead Code
Insertion forces the model to decide which inserted statements are semantically
inert and removable. Literal Encryption forces recovery of hidden values by
reasoning over the inserted decryption logic. The three tasks therefore span
lexical, structural, and semantic inverses while the forward tool is held
fixed.

\subsection{Additional Settings}
\label{sec:additional_settings}

Two further checks operate outside obfuscation. For translation, we fine-tune
GPT-4.1-mini on 1{,}000 Python-to-Java pairs drawn from matched CodeNet
Python800/Java250 problems and evaluate the untrained reverse direction on
300 ConDefects programs~\cite{condefects}, counting a program as solved when at least
half of its tests pass. For compilation, we fine-tune the same model on 1{,}000
C-to-assembly pairs from decompile-bench and evaluate the untrained reverse
direction on 164 HumanEval-Decompile functions through re-execution against unit
tests. Neither check repeats the full multi-model BRF protocol. Each one tests
whether the directional pattern survives outside the main Java obfuscation
study.

\subsection{Datasets}
\label{sec:datasets}

We need executable programs with reliable test suites because semantic
equivalence is the only defensible ground truth for both obfuscation and
deobfuscation. We therefore use CodeNet~\cite{puri2021codenet} for training and
ConDefects~\cite{condefects} for evaluation. From CodeNet we select
20{,}000 Java programs through stratified sampling and keep only programs that
compile successfully and have usable tests. 
From ConDefects we use 300 Java
programs, each with an average of 16 tests. Every evaluation program is passed
through all three obfuscation types.

The split between training and evaluation data is deliberate. CodeNet supplies
breadth for fine-tuning, and ConDefects supplies held-out programs for
validation. The design lets us ask two questions at once, whether a model learns
a forward transformation that transfers to new programs, and whether the same
adaptation preserves any inverse capability on unseen code. The first answer is
frequently affirmative. The second answer is the central finding of the paper.

\subsection{Fine-Tuning Protocol}
\label{sec:training_protocol}

For the open-source models we use LoRA~\cite{hu2022lora}. We avoid full
fine-tuning, which is far more costly at this scale and forgets more of the base
model~\cite{biderman2024lora}, because the paper concerns a post-training risk
that practitioners meet under common adaptation regimes rather than under the
most expensive one. We avoid prefix tuning~\cite{li2021prefix} as well, since
the mechanism question concerns the generation policy after adaptation, and LoRA
alters the parameters used during generation. We also tried QLoRA~\cite{dettmers2023qlora}, but the fine-tuned models often produced degenerate outputs made of incoherent characters, so we did not retain it.
This observation is consistent with prior work reporting that quantization can degrade LoRA fine-tuning~\cite{li2024loftq, qin2024irqlora}.

We calibrated the training size and schedule in a preliminary study. For the
open-source models, we compared 5{,}000, 10{,}000, and 20{,}000 forward pairs per
transformation, and forward performance was highest with 10{,}000 pairs. Training
loss and held-out forward accuracy saturated by the third epoch, so we train for
three epochs. For the proprietary API models, we run
instruction fine-tuning through the provider interface~\cite{openai_sft_guide}.
The provider reports gains from as few as 50 to 100 examples and recommends
scaling the dataset according to observed results. Following this guidance, we
compared 500, 1{,}000, and 2{,}000 pairs, and forward performance was highest
with 1{,}000 pairs. The resulting size asymmetry between open-source and
proprietary runs reflects these calibrations. It limits direct comparison across
model families, but not the central claim, since the directional failure appears
in both. Decoding always uses temperature zero, so that performance reflects the
adapted model and not sampling noise.

We compare three training strategies. Standard fine-tuning (SFT) trains on the
forward mapping from original to obfuscated code alone. Balanced bidirectional
fine-tuning (BFT) adds explicit supervision in both directions and so establishes
a recovery ceiling for the reverse task. Contrastive fine-tuning (CFT) keeps the
forward generation objective and adds semantic triplets
$(C_{\text{anchor}}, C_{\text{positive}}, C_{\text{negative}})$, where the
positive is an obfuscated version of the anchor and the negative is an unrelated
program. CFT is not designed to maximize forward accuracy. It tests whether
preserving semantic proximity in representation space reduces the reverse failure
without any reverse labels.
\subsection{Hyperparameters}
\label{sec:hyperparameters}

Table~\ref{tab:hyperparameters} lists the main hyperparameters. Open-source
models use LoRA rank 16 and $\alpha=32$, a configuration standard for
7B-scale adaptation~\cite{hu2022lora,dettmers2023qlora}. All open-source runs
use learning rate $2\times10^{-4}$, batch size 4, and three epochs. API-based
models use provider-controlled optimization with the listed exposed settings.
For CFT, the contrastive and generation losses are equally weighted on the
open-source runs.

\begin{table}[tb]
\centering
\caption{Hyperparameter configuration for all models and training strategies.}
\label{tab:hyperparameters}
\begin{adjustbox}{width=.6\columnwidth,center}
\begin{tabular}{@{}lcccccc@{}}
\toprule
\textbf{Model} & \textbf{LoRA $r$} & \textbf{LoRA $\alpha$} & \textbf{LR} & \textbf{Batch} & \textbf{Epochs} & \textbf{Temp.} \\ \midrule
Qwen2.5-7B & 16 & 32 & 2e-4 & 4 & 3 & 0 \\
Qwen2.5-Coder-7B & 16 & 32 & 2e-4 & 4 & 3 & 0 \\
DeepSeek-R1-Distill-Qwen-7B & 16 & 32 & 2e-4 & 4 & 3 & 0 \\
Mistral-7B & 16 & 32 & 2e-4 & 4 & 3 & 0 \\
GPT-3.5-Turbo & --- & --- & auto & 4 & 3 & 0 \\
GPT-4.1-mini & --- & --- & auto & 4 & 3 & 0 \\ \bottomrule
\end{tabular}
\end{adjustbox}
\vspace{-0.1cm}
\end{table}

\subsection{Evaluation Metrics}
\label{sec:metrics}

No single metric shows whether code has been obfuscated or deobfuscated.
We therefore combine a semantic check with several complementary syntactic measures.

\begin{itemize}[leftmargin=*]

\item \textbf{Semantic correctness.}
An output is valid if it \wcircle{1} compiles, \wcircle{2} runs without runtime errors, and \wcircle{3} passes all test cases of the original program.
This check applies to both directions and confirms that the transformation preserves behavior.

\item \textbf{Syntactic similarity ($S$).}
CodeBLEU~\cite{ren2020codebleu} measures token-level and structural similarity between two programs.
For obfuscation, a lower $S(C_{\text{obf}}, C_{\text{orig}})$ indicates a stronger transformation.
For deobfuscation, a higher $S(C_{\text{deobf}}, C_{\text{orig}})$ indicates better recovery.
Here, $C_{\text{obf}}$ and $C_{\text{deobf}}$ denote the model outputs for each direction.

\item \textbf{Readability ($R$).}
The model of Scalabrino et al.~\cite{scalabrino2018readability} estimates how easily a human can read a program, with a score in $[0,1]$.
It requires no reference program.
A successful deobfuscation should raise $R(C_{\text{deobf}})$ above $R(C_{\text{obf}})$ and bring it closer to $R(C_{\text{orig}})$.

\item \textbf{AST-based measures.}
We parse programs with tree-sitter~\cite{tree-sitter} and compare sets of AST elements with the Jaccard index, $J(A,B) = |A \cap B| / |A \cup B|$.
The two measures use opposite directions.
\begin{itemize}
  \item $AST_{\text{synt}} = 1 - J(I_A, I_B)$ compares the sets of identifiers $I$, such as variable and method names.
  It measures lexical divergence: 0 means identical identifiers, and 1 means that all identifiers were replaced.
  \item $AST_{\text{struct}} = J(N_A, N_B)$ compares the sets of control-flow nodes $N$, such as loops, conditionals, and declarations.
  It measures structural similarity: 1 means an identical skeleton, and 0 means a fully different one.
\end{itemize}
A higher $AST_{\text{synt}}$ therefore signals more renaming, whereas a lower $AST_{\text{struct}}$ signals more structural change.

\item \textbf{Complexity.}
We also report token count, computed with a BERT-base tokenizer~\cite{devlin2018bert}, and cyclomatic complexity, computed with Lizard~\cite{lizardtool}.
Obfuscation tends to increase both, and deobfuscation should reduce them.
% These two measures are descriptive and do not enter the decision rule.

% \end{itemize}


% Together, these metrics allow us to evaluate transformation quality from multiple angles, semantic fidelity, syntactic distortion, structural recovery, and human interpretability, providing a framework for assessing both forward and reverse reasoning capabilities.

These two measures are descriptive and do not enter the reverse decision rule.

\end{itemize}

\textbf{Reverse recovery decision rule.}
Semantic validity alone is not enough.
An unchanged obfuscated input also passes the same tests.
Reverse success requires the semantic check plus joint movement toward the original on all three channels: CodeBLEU, identifiers, and structure.
We require $S(C_{\text{deobf}}, C_{\text{orig}}) > S(C_{\text{tool}}, C_{\text{orig}})$, $AST_{\text{synt}}(C_{\text{deobf}}, C_{\text{orig}}) < AST_{\text{synt}}(C_{\text{tool}}, C_{\text{orig}})$, and $AST_{\text{struct}}(C_{\text{deobf}}, C_{\text{orig}}) > AST_{\text{struct}}(C_{\text{tool}}, C_{\text{orig}})$.
The comparison is per program against that program's tool output.
Reported mean tool CodeBLEU values are $0.50$ (Identifier Renaming),
$0.33$ (Dead Code Insertion), and $0.22$ (Literal Encryption)
(Table~\ref{tab:reversibility_heatmap_data}).

Together, these metrics allow us to evaluate transformation quality from multiple angles, semantic fidelity, syntactic distortion, structural recovery, and human interpretability, providing a framework for assessing both forward and reverse reasoning capabilities.